\section{Three problems}
\label{sec:problem}

We state three problems. Problem 1 is that one literature reaches incompatible
conclusions about a single attack. Problem 2 is that two literatures define the
same quantity and neither cites the other. Problem 3 is that a stated threat model
does not constrain the attack it was written for. Section~\ref{sec:problem-cause}
takes the three in turn and shows that one absence produces all of them.

\subsection{Problem 1: incompatible conclusions within one literature}
\label{sec:problem-one}

Membership inference decides whether one record was in a model's training set.
Meeus et al.\ count at least \MIAnewMethods{} new methods for it proposed in 2024
alone \cite{MSJFRM25}, and the four verdicts quoted in Section~\ref{sec:intro} span
the whole range from success to unfalsifiability. What produced that spread is
visible once each paper's adversary is written down.

Carlini et al.\ give their adversary data drawn from the target's distribution and
the budget to train many shadow models on it, then compare the target's per-example
loss against the two resulting distributions \cite{CCNSTT22}. Duan et al.\ give
theirs a benchmark whose members and non-members come from the Pile's own train and
test split, which removes the shortcut that a benchmark assembled after training
leaves open \cite{DSMMSZTCEH24}. Das et al.\ give theirs no model access at all,
and it still wins, which places the signal in the benchmark rather than in the
model \cite{DZT25}. Meeus et al.\ locate that signal precisely, reporting
\MIAwikiLo{} to \MIAwikiHi{} AUC on WikiMIA from a bag-of-words classifier whose
most predictive features are years \cite{MSJFRM25}.

Zhang et al.\ reach the strongest verdict \cite{ZDKT25}. To bound a false positive
rate, the party running the attack must know the distribution of its test statistic
under the hypothesis that the model was trained without the record in question.
That hypothesis cannot be sampled.

None of the five states a budget in a form the others could be checked
against. The last of them comes closest to stating one. Zhang et al., in a position paper
not yet refereed, give two independent
reasons the null hypothesis cannot be sampled, and the second survives the removal
of the first. Granted full knowledge of the training set and of the training
algorithm, sampling the null ``would still be prohibitively expensive as it would
require training new large models'' \cite{ZDKT25}. That is a security claim
declared unfalsifiable against a party that cannot afford a training run. The paper
names the cost and moves on. It does not record the bound as a parameter, so the
claim cannot be compared with a claim made against a party of any other budget.

\subsection{Problem 2: the attacker-to-honest cost ratio, defined twice without contact}
\label{sec:problem-two}

A proof-of-learning asks a party claiming to have trained a model to produce a
record of the training run, which a verifier then checks \cite{JYCDTCP21}. The
scheme is secure when forging the record costs at least as much as training
honestly. Jia et al.\ state that requirement three times, in the form
$\mathbb{E}[C_A] \geq \mathbb{E}[C_T]$, where $C_A$ is the cost of any spoofing
strategy and $C_T$ the cost of the honest training run.

The inequality carries no exponent, no adversary class, no security parameter and
no assumption. As written it holds by definition for the spoof that simply retrains the
model, and the authors say plainly that they cannot prove it for anything else.
Having concluded in \S VII-B3a that a passing forgery ``would be at least as
computationally expensive as what it took'' the honest trainer, they add two pages
later that ``given that we are not aware of a mechanism to prove this formally, we
present the following as a conjecture'' \cite{JYCDTCP21}.

Fine-grained cryptography defines the same quantity and gives it a name. Brzuska
and Couteau's Definition 6 fixes an $(\varepsilon,\delta)$-fine-grained one-way
function as one that no circuit of size below $|f|^{1+\delta}$ inverts with
non-negligible probability \cite{BC22}. LaVigne, Lincoln and Vassilevska Williams
name the exponent itself, in their Definition 5, as ``the (relative) gap of an
$T(n)$ fine-grained one-way function $f$'' \cite{LLW19}. Around that one constant
they build fine-grained key exchange and fine-grained public-key encryption.

The two literatures are describing the same object. Neither cites the other. The
proof-of-learning line has the quantity in hand and no name for it, so it cannot
say how much more expensive forgery is, and the fine-grained line has the name and
no application to learning, so nothing in it has ever been instantiated on a
training run.

\subsection{Problem 3: threat models stated without a resource bound}
\label{sec:problem-three}

Athalye et al.\ take the \OGpopulation{} non-certified defences accepted at ICLR
2018 that claim white-box robustness \cite{ACW18}. \OGobfuscated{} of them rely on
a phenomenon the authors name obfuscated gradients. \OGbrokenFull{} are
circumvented completely and \OGbrokenPart{} partially, each under the threat model
its own paper stated.

Two years later Tram\`er et al.\ repeat the exercise on \AAdefences{} defences from
ICLR, ICML and NeurIPS, all of which had performed adaptive evaluations
\cite{TCBM20}. They circumvent all \AAdefences{}. Their diagnosis is that they needed nothing
new. In their words, ``the technical tools to evaluate defenses properly
already exist'', and the issue ``is methodological rather than technical''
\cite{TCBM20}.

A defence and its break can both be correct. The threat models in question fix the
perturbation set and the access level, and they say nothing about what the attacker
may compute or about what counts as success. A model that fixes neither of those
does not constrain the attack, so a defence can satisfy it and an attack can
violate the defence without either party making an error.

\subsection{The common cause}
\label{sec:problem-cause}

The three problems share one cause. Most of these works state no budget for any
party, so there is nothing to hold a view about. Researchers holding different
views about what an adversary is would already be an improvement, because such
positions can be compared. Section~\ref{sec:findings} reports how many.

Take the three in turn. In Problem 1 the five membership inference results grant
their adversaries different things, and each verdict follows its grant. Carlini et
al.\ grant the target's distribution and the compute to train shadow models on it.
Duan et al.\ grant a benchmark that removes the shortcut. Das et al.\ grant no
model access at all. The results cannot be ordered, because no two of them price
the same party in the same unit.

In Problem 2 the proof-of-learning line holds the quantity and has no name for it,
while the fine-grained line holds the name and has never instantiated it on a
training run. Neither can find the other by searching, because neither shares
vocabulary with the other. The shared object is a ratio between two parties' work,
and only one of the two literatures writes it as an exponent.

In Problem 3 the threat models fix the perturbation set and the level of access,
and they leave open what the attacker may compute and what counts as success. A
defence can satisfy such a model and an attack can then break the defence, with
neither party in error, because the model never ruled the attack out.

The field has a name for the third case and not for the first two. Arp et al.'s
tenth pitfall, quoted in Section~\ref{sec:intro}, is the threat model, defined as
the case where ``the security of machine learning is not considered, exposing the
system to a variety of attacks, such as poisoning and evasion attacks''
\cite{AQPWPWCR22}. Neither that pitfall nor its recommendations quantify the
adversary by computation, memory, samples or queries. The defect the field
diagnoses in its own threat models is that they are unstated or imprecise, and
never that they are unpriced.
